\documentclass[a4paper]{article}
\usepackage{amsfonts}
\usepackage{amsmath}
\usepackage{amssymb}
\usepackage{graphicx}     

\begin{document}
  
\noindent \large Auteur: Eryk Borczuch \\
\noindent \large Studierichting: Informatica\\
\noindent \large Oefeningenreeks 1D nr. 19.5\\

\medskip

\normalsize

$z^3 - z^2 - z -2$\\

\begin{tabular}{c|ccccccccc}
& $1$ & $-1$ & $-1$   & \vline & $-2$ \\
$2$ & & $2$ & $2$ & \vline & $2$ \\ \hline
& $1$ & $1$ & $1$ &    \vline & $0$ \\
\end{tabular}\\

$= (z^2+z+1)(z-2)$\\

nulpunten: \\

$z_1 = 2$\\

$D = 1^2 -4\cdot1\cdot1$\\

$= -3$\\

$z_{2,3} = \dfrac{-1 \pm \sqrt{3}i}{2}$

dus de nulpunten zijn: $\left(2, \dfrac{-1 + \sqrt{3}i}{2}, \dfrac{-1 - \sqrt{3}i}{2}\right)$
\begin{thebibliography}{999}
%\bibitem{a} Referentie
\end{thebibliography}
\end{document} 
